\section{Open problems}\label{sec:open}

\paragraph{$\PR$ versus $\NPR$.} If a language with short real certificates is the graph of a polynomial of polynomial degree, or the region on one side of such a graph, then \cref{thm:graph,thm:threshold} give an upper bound on the border complexity of that polynomial. Therefore a lower bound on border complexity cannot prove that such a language lies outside $\PR$. If a tree of depth $t$ decides a language whose boundary contains an irreducible hypersurface $\{h=0\}$, then the generic-path argument gives a nonzero multiple of $h$ that has a straight-line program of length $O(t)$, but with no bound on its degree \cite{Bur04}.

B\"urgisser conjectures that $\VNP\not\subseteq\overline{\VP}$ over $\CC$ implies $\PC\ne\NPC$ nonuniformly \cite[Conjecture~1.1]{Bur26}. The hypothesis follows from \cref{thm:lowerbound} in its form over $\CC$. Polynomial-size circuits for the resultant of $n+1$ quadratic forms imply $\VP=\VNP$ over $\CC$ \cite{Bur26}. Extending this implication to every nonzero multiple of that resultant would prove the conjecture. The resultant has exponential degree, so the degree dependence in \cref{thm:burgisser} is an obstacle. See also \cite{SS95} for the analogous role of multiples of integers.

\paragraph{B\"urgisser's conjecture and squarefree parts.} We record a side observation: B\"urgisser's conjecture follows from a bound on the circuit complexity of squarefree parts. Unlike the rest of the paper, this observation is not formalized in Lean. For a nonzero polynomial $P$ over $\CC$ write $\rad(P)$ for the product of its distinct irreducible factors, and write $\doteq$ for equality up to a nonzero scalar. Consider the following hypothesis.
\begin{enumerate}[label=\textup{(R)},leftmargin=3em]
\item There is a polynomial $\pi$ such that, if a polynomial $P$ over $\CC$ is computed by an arithmetic circuit of size $s$, then $\rad(P)$ is computed by an arithmetic circuit of size $\pi(s)$, whatever the degree of $P$.
\end{enumerate}
Let $\overline{\mathrm{(R)}}$ be the weaker hypothesis that bounds the border complexity of $\rad(P)$ by $\pi(s)$. Dutta, Saxena and Sinhababu conjecture a polynomial bound for the complexity of the squarefree part \cite[Conjecture~1]{DSS22}. In a form that bounds the minimum of the degree and the circuit size of $\rad(P)$, it implies \textup{(R)}: the polynomial $P\cdot(Z^{2^s}-1)$, with a new variable $Z$, has a circuit of size $O(s)$ and a squarefree part of degree at least $2^s$, so for large $s$ the bound applies to its circuit size, and setting $Z=0$ gives $\rad(P)$ up to a scalar.

\begin{lemma}\label{lem:radpath}
Let $f\in\CC[x]$ be irreducible, and let a complex tree of depth $t$ decide $\{f=0\}$. Then there are nonzero polynomials $A,B,C\in\CC[x]$, computed together by one straight-line program with $O(t)$ operations, such that
\[
f\ \doteq\ \frac{\rad(AB)\,\rad(CB)}{\rad(ACB)\,\rad(B)} .
\]
\end{lemma}

\begin{proof}
Run the tree at the generic point of $V(f)$, that is, over the function field of the irreducible hypersurface $V(f)$, and represent the values on the path it follows by fractions as in \cref{sec:cells}. The points of a dense open subset of $V(f)$ follow the same path, so the path accepts. Its conditions are $E_1=\dots=E_m=0$ and $B\ne0$, where the $E_i$ are the tested numerators of the tests answered ``yes'' and $B$ is the product of the other tested numerators and of the numerators of the divisors. All of them are computed by one straight-line program with $O(t)$ operations, $f$ divides every $E_i$, and $f$ does not divide $B$. Every point that satisfies the conditions is accepted, so it lies in $V(f)$. We discard the $E_i$ that are the zero polynomial. Some $E_i$ remains, since otherwise the open set $\{B\ne0\}$ would be accepted.

Let $A$ be one of the remaining $E_i$. Let $C$ be a nonzero linear combination of the $E_i$ that is not divisible by any irreducible factor of $A$ that fails to divide some $E_j$. Each such factor excludes a proper linear subspace of coefficient vectors, and there are finitely many. The common irreducible factors of $A$ and $C$ are then exactly the common irreducible factors $g$ of all the $E_i$. If such a $g$ does not divide $B$, then a dense open subset of $V(g)$ satisfies the conditions of the path, so $V(g)\subseteq V(f)$ and $g\doteq f$. Conversely $f$ divides $A$ and $C$ and does not divide $B$. For an irreducible $g$ the exponent of $g$ in the displayed quotient is
\[
\mathbf 1_{g\mid AB}+\mathbf 1_{g\mid CB}-\mathbf 1_{g\mid ACB}-\mathbf 1_{g\mid B},
\]
which is $1$ if $g$ divides $A$ and $C$ but not $B$, and $0$ otherwise.
\end{proof}

\begin{proposition}\label{prop:radical}
Assume \textup{(R)}. If $\NPC$ is contained in nonuniform polynomial time over $\CC$, then $\VP=\VNP$ over $\CC$. Under $\overline{\mathrm{(R)}}$ the same assumption gives $\VNP\subseteq\overline{\VP}$ over $\CC$. In particular $\overline{\mathrm{(R)}}$ implies B\"urgisser's conjecture \cite[Conjecture~1.1]{Bur26}, and together with \cref{thm:lowerbound} in its form over $\CC$ it implies $\PC\ne\NPC$ nonuniformly.
\end{proposition}

\begin{proof}
Let $p(T')$ be a monic polynomial of degree $n$ with roots $r_1,\dots,r_n$, let $q_1,\dots,q_n$ have degree less than $n$, and let $T$ be a further variable. Put
\[
f(p,q,T)=\prod_{\sigma\in S_n}\Bigl(1-T\prod_{i=1}^nq_i(r_{\sigma(i)})\Bigr).
\]
Its coefficients are symmetric in the $r_j$, so $f$ is a polynomial in $T$ and in the coefficients of $p$ and of the $q_i$. Treat the $r_j$ and the coefficients of the $q_i$ as independent variables. The values $q_i(r_j)$ are then independent over $\CC(r)$ by the Vandermonde determinant, so the $n!$ products $\prod_iq_i(r_{\sigma(i)})$ are distinct, and $S_n$ permutes them transitively. Their common minimal polynomial over the field generated by the coefficients of $p$ and of the $q_i$ therefore has degree $n!$. It equals $T^{n!}f(1/T)=\prod_\sigma\bigl(T-\prod_iq_i(r_{\sigma(i)})\bigr)$, which is monic in $T$ and hence irreducible as a polynomial by Gauss's lemma. Since $f(0)=1$, also $f$ is irreducible. The set $\{f=0\}$ is in $\NPC$: guess $z_1,\dots,z_n$, check the coefficient identity $p(T')=\prod_j(T'-z_j)$, and check $T\prod_iq_i(z_i)=1$.

By assumption, complex trees of depth polynomial in $n$ decide $\{f=0\}$. \Cref{lem:radpath} and \textup{(R)} give polynomial-size circuits for $U=\rad(AB)\rad(CB)$ and $V=\rad(ACB)\rad(B)$, with $U=\lambda fV$ for some $\lambda\ne0$. Each radical is divisible at most once by $T$, so $V$ has order $a\le2$ in $T$ at $0$. Write $U_j$ and $V_j$ for the coefficients of $T^j$. Since $f(0)=1$ we have $U_a=\lambda V_a\ne0$ and
\[
[T]\,f=\frac{U_{a+1}}{U_a}-\frac{V_{a+1}}{V_a}.
\]
Fix distinct complex numbers $r_1,\dots,r_n$ for which $V_a$ is a nonzero polynomial in the coefficients of the $q_i$, and determine the $q_i$ by interpolation from $q_i(r_j)=X_{ij}$. Then $[T]f=-\sum_\sigma\prod_iX_{i\sigma(i)}=-\per_n(X)$. The coefficients of $T^j$ for $j\le3$ are computed from the circuits for $U$ and $V$ in $\CC[x][T]/(T^4)$. Choose $X_0$ with $V_a(X_0)\ne0$ and evaluate the right-hand side at $X_0+\tau W$ modulo $\tau^{n+1}$; the constant terms of the denominators are nonzero constants, so the divisions cost $O(n^2)$ operations each. Summing the coefficients gives $-\per_n(X_0+W)$ exactly, a polynomial-size circuit for the permanent. As the permanent is complete for $\VNP$, $\VP=\VNP$. Under $\overline{\mathrm{(R)}}$ the same construction applied to approximating circuits gives coefficients that converge to those of the permanent, because the denominators have nonzero limits at $X_0$. Hence the permanent has polynomial border complexity and $\VNP\subseteq\overline{\VP}$.
\end{proof}

The identity of \cref{lem:radpath} holds unconditionally. Hypothesis \textup{(R)} is used only to bound the coefficients of $T^j$, $j\le3$, of the four squarefree parts on the interpolation slice. The identity depends on discarding multiplicities: with the polynomials $AB$, $CB$, $ACB$, $B$ in place of their squarefree parts the quotient is $1$, and quantities that are additive under products, such as root sums counted with multiplicity, cancel in the same way.

Lower bounds for all nonzero elements are known for other ideals. Forbes, Shpilka, Tzameret and Wigderson \cite{FSTW16} prove lower bounds for all nonzero multiples of explicit polynomials in several restricted circuit models. Andrews and Forbes \cite{AF22} show that every nonzero polynomial in the ideal of $r\times r$ minors of a generic matrix yields, through a small oracle circuit of depth three, an approximation of a determinant of size $\Theta(r^{1/3})$. For the principal ideal $(Y-\per_m(X))$ over $\RR$, \cref{thm:lowerbound,thm:burgisser} give a lower bound even for unrestricted arithmetic circuits. If a nonzero multiple has a straight-line program of length $\ell$, then the border complexity of $\per_m$ is at most $Cm^3(\ell+m^2)$. Hence $\ell\ge 2^{\Omega((\log m)^{4/3})}$ regardless of the degree of the multiple.

\paragraph{The universal--existential level.} Suppose $y=f(x)$ holds exactly when for all $z$ some $u$ is accepted. Off the graph there is a $z$ for which every $u$ is rejected. Such witnesses can escape to infinity as $(x,y)$ approaches the graph. For example, let the verifier test $yz-1=0$, accepting exactly when this test fails and ignoring $u$. Then $\forall z\,\exists u$ defines $y=0$. For $y\ne0$ the only rejecting value is $z=1/y$. On any fixed closed ball $B$ in $z$, the accepting margin of \eqref{eq:margin} is $v(y,z)=\min\{1,(1-yz)^2\}$, which tends uniformly to $1$ as $y\to0$. For fixed $N$ and $a$, the objective \eqref{eq:J} at maximizing margins changes from its value at $y=0$ by at most $\sup_{z\in B}|v(y,z)^N-1|$, uniformly over all atom locations and weights in that ball. Its maximum is therefore continuous across $y=0$. It remains open whether the border complexity of $f$ is polynomially bounded in its degree, the number of variables, the tree depth and the two block lengths for this quantifier order.

\paragraph{Thresholds at the second level.} For the permanent, negating one row reduces either closed threshold to the other, and \cref{cor:second} gives the lower bound for both. For a general polynomial $f$, does an $\exists z\,\forall u$ verifier for $\{y\le f(x)\}$ yield the same polynomial bound on border complexity as in \cref{thm:second}? Here the pairs $(y,z)$ over $x$ for which every $u$ is accepted can form a positive-dimensional set. If $v^N$ is bounded by a multiple of the squared distance to this set, the envelope vanishes exactly on it. One possible approach is to minimize $e-\vartheta y$ with a smaller generic linear perturbation to select a boundary point.

\paragraph{Higher levels.} At the third level, the corresponding set of lower dimension is the set of points at which an accepting margin is positive for every value of an intermediate block of variables. Its complement is not described by a single maximization. It is open whether the graph of the permanent lies in $\PHR$. If it does not, then $\PHR\subsetneq\PARR$. It is also open whether $\PHR$ is contained in polynomial time with an oracle for the permanent; such an inclusion would be a real analogue of Toda's theorem. Compare \cite{BM16,Bur26}.

\paragraph{Rounds.} \Cref{thm:amper} gives a lower bound of $\Omega(\log\log m)$ on the number of rounds of a public-coin protocol for the permanent, while the upper bound is $m$. The lower bound is limited by the growth of the sample sizes in \cref{lem:coins}. We do not know whether a constant number of rounds with private coins and real messages can be reduced to a real certificate.
